
\section{Estimating the Curvature Factors}
\label{app:matnormal-t4}
\Cref{alg:ours} turns the single-batch update of
\Cref{sec:sample-loss-lmo} into an optimizer by averaging its two inputs,
the gradient in the linear objective and the second moment in the
outergradient constraint. Averaging over recent steps draws on more samples
while keeping the estimate local to the current iterate. The factor-gradient
momenta $\widehat M_A,\widehat M_B$ supply
the projections used in the direction step~\eqref{eq:factor-directions},
while an EMA accumulates the vectors $p,q$ that define the diagonal
curvature factors $P,Q$. \Cref{sec:weighted-estimation-t4} derives these
quantities from LoRA factor gradients alone, and
\Cref{sec:scale-placement-t4} explains why the accumulated vectors are
normalized after averaging. The update is invariant to the scale of the
curvature factors, but the average that estimates them depends on that scale.

\subsection{Estimation from the Factor Gradients}
\label{sec:weighted-estimation-t4}

\paragraph{Running averages.} Let $s$ index optimizer steps, let $G_s$ be
the batch gradient at step $s$, and write $g_s=\vec(G_s)$.
\Cref{alg:ours} uses look-ahead momentum for the gradient input, matching
line~\ref{ln:mom},
\[
  M_{t+1}=\beta_1M_t+(1-\beta_1)G_t,
  \qquad
  \widehat M_t=\beta_1M_{t+1}+(1-\beta_1)G_t.
\]
For the curvature input it uses an exponential moving average, which keeps
memory constant and downweights older gradients computed farther from the
current iterate. We write the induced weights on past steps as
\[
  \EMA_{s \leq t}[\phi(G_s)]=\sum_{s\le t}w_{t,s}\,\phi(G_s),
  \qquad
  w_{t,s}=(1-\beta_2)\beta_2^{t-s}.
\]
Thus the curvature factors are fit to the average
\[
  \overline{\Sigma}_t=\EMA_{s\leq t}[g_sg_s^\top],
\]
rather than to the current batch alone.

\paragraph{Using the factor gradients.} LoRA backpropagation gives
$G_{A,s}=B_s^\top G_s$ and
$G_{B,s}=G_sA_s^\top$~\eqref{e-grad-wrt-factors}. Forming the dense batch
gradient $G_s$ would require a $d_{\mathrm{out}}\times d_{\mathrm{in}}$
matrix per layer, so the optimizer builds the averages from the factor
gradients. The single-batch derivation approximates the per-sample second
moment $\Sigma$ by $Q\otimes P$~\eqref{eq:pq-factor}. The optimizer applies
the same model to $\overline{\Sigma}_t$, and three approximations make its
factors estimable from $G_{A,s}$ and $G_{B,s}$.
\begin{enumerate}[leftmargin=1.6em,itemsep=2pt,topsep=2pt]
  \item Let
  $g_{s,i}=\vec(\nabla\ell_{x_{s,i}}(W_s))$ be the per-sample gradients in
  the size-$n_s$ minibatch at step $s$. Replace their second moment by the
  outer product of the batch gradient,
  \[
    \frac1{n_s}\sum_{i=1}^{n_s}g_{s,i}g_{s,i}^\top
    \;\leadsto\;
    g_sg_s^\top ,
  \]
  the batch-gradient second-moment approximation used by
  AdaGrad~\cite{adagrad}, Adam~\cite{adam}, and Shampoo~\cite{shampoo}.
  It is exact at batch size one, and the accumulated second moment is then
  $\overline{\Sigma}_t$.
  \item Model $\overline{\Sigma}_t$ by diagonal Kronecker factors. We write
  $P^\star,Q^\star$ for these model factors to distinguish them from the
  estimates $P,Q$ in \Cref{alg:ours}. Since
  $Q^\star\otimes P^\star$ is unchanged by the
  rescaling $(P^\star,Q^\star)\mapsto(aP^\star,a^{-1}Q^\star)$, we fix
  $\norm{\diag P^\star}_\infty=\norm{\diag Q^\star}_\infty=1$ and use a
  scalar $\kappa_t>0$ for the magnitude,
  \begin{equation}\label{e-diag-kron-second-mom-model-t4}
    \EMA_{s \leq t}[g_sg_s^\top]\approx\kappa_t\,(Q^\star_t\otimes P^\star_t).
  \end{equation}
  \item When relating moments of $G_{A,s}$ and $G_{B,s}$ to $\overline{\Sigma}_t$, replace
  the past LoRA factors by the current factors inside the weighted averages,
  \[
    A_s\;\leadsto\;A_t,\qquad B_s\;\leadsto\;B_t,\qquad s\le t,
  \]
  which is accurate when the factors change slowly over the EMA window.
\end{enumerate}

\paragraph{Factor momenta.} The objective decomposes, under the linearized
product update~\eqref{eq:DeltaW-lin} with $\widehat M_t$ replacing the batch
gradient in~\eqref{eq:polorafull}, as
\[
  \ip{\widehat M_t,\,\Delta B\,A+B\,\Delta A}
  =\ip{\widehat M_tA^\top,\Delta B}+\ip{B^\top\widehat M_t,\Delta A},
\]
so the update depends on $\widehat M_t$ only through $\widehat M_tA^\top$
and $B^\top\widehat M_t$. The momenta $\widehat M_A,\widehat M_B$ in
\Cref{alg:ours} apply the same look-ahead rule to $B_s^\top G_s$ and
$G_sA_s^\top$. Replacing $A_s,B_s$ by $A_t,B_t$ in the weighted sums gives
$\widehat M_A\approx B_t^\top\widehat M_t$ and
$\widehat M_B\approx\widehat M_tA_t^\top$.

\paragraph{Moments of the factor gradients.} To estimate $P^\star$ and
$Q^\star$ from factor gradients, we link their second moments to $\overline{\Sigma}_t$.
Since
$\vec(G_{A,s})=(I\otimes B_s^\top)\vec(G_s)$,
\begin{align*}
 \Sigma_{A,t} &:=  \EMA_{s \leq t}\big[ \vec(G_{A,s})  \vec(G_{A,s})^\top\big]  \\
 &=  \EMA_{s \leq t}\big[ (I\otimes B_s^\top)\vec(G_s)\vec(G_s)^\top   (I\otimes B_s)\big]  \\
   &\approx  (I\otimes B_t^\top) \EMA_{s \leq t}\big[ g_s g_s^\top   \big]  (I\otimes B_t),
\end{align*}
where the last line uses $B_t$ in place of $B_s$ inside the weighted
average. Applying the diagonal Kronecker model~\eqref{e-diag-kron-second-mom-model-t4}
then gives
\begin{align*}
 \Sigma_{A,t}
   &\approx
   \kappa_t\,
   (I\otimes B_t^\top)
   (Q^\star_t\otimes P^\star_t)
   (I\otimes B_t) \\
   &= \kappa_t\,Q^\star_t\otimes(B_t^\top P^\star_tB_t),
\end{align*}
where the equality is the mixed-product rule of Kronecker products. The
analogous calculation for $B$ gives
\begin{align}
  \Sigma_{A,t}
    &\approx \kappa_t\,Q^\star_t\otimes C^\star_{A,t},
  \nonumber \\
  \Sigma_{B,t}
    &:=
      \EMA_{s \leq t}\big[
        \vec(G_{B,s}^\top)\vec(G_{B,s}^\top)^\top
      \big]
    \approx \kappa_t\,P^\star_t\otimes C^\star_{B,t},
  \label{eq:observable-moments-t4}
\end{align}
where $C^\star_{A,t}=B_t^\top P^\star_tB_t$ and
$C^\star_{B,t}=A_tQ^\star_tA_t^\top$ are the $r\times r$ matrices
of~\eqref{eq:ours-closure} evaluated at the model factors.

\paragraph{Fitting the factors.} We fit $P^\star$ and $Q^\star$ to the
moments in~\eqref{eq:observable-moments-t4}. In this paragraph all
quantities are at the current step, so we suppress the subscript $t$.
Following KL-Shampoo~\cite{klshampoo}, we use the log-determinant
divergence~\cite{logdetdiv}
\[
  D(\Sigma,S)=\Tr(\Sigma S^{-1})-\log\det(\Sigma S^{-1})-d,
\]
defined for positive definite matrices $\Sigma$ and $S$ of common dimension
$d$. The $A$-gradient moment identifies $Q^\star$ once
$C^\star_A=B^\top P^\star B$ is fixed, and the $B$-gradient moment
identifies $P^\star$ once $C^\star_B=AQ^\star A^\top$ is fixed. Each fit
therefore freezes the opposite factor at its current estimate,
\begin{equation}\label{eq:factor-fit-t4}
  \hat{q}(\widehat P)
  =\argmin_{q>0}\;
  D\big(\Sigma_A,\; \diag(q) \otimes \widehat C_A\big),
  \qquad
  \hat{p}(\widehat Q)
  =\argmin_{p>0}\;
  D\big(\Sigma_B,\;\diag(p) \otimes \widehat C_B\big),
\end{equation}
with $\widehat C_A=B^\top\widehat PB$ and $\widehat C_B=A\widehat QA^\top$
for diagonal $\widehat P,\widehat Q\succ0$. The blockwise first-order
conditions give the closed forms
\[
  \hat{q}(\widehat P)=\frac1r\diag\!\left(\EMA_{s \leq t}[G_{A,s}^\top \widehat C_A^{-1}G_{A,s}]\right),
  \qquad
  \hat{p}(\widehat Q)=\frac1r\diag\!\left(\EMA_{s \leq t}[G_{B,s} \widehat C_B^{-1}G_{B,s}^\top]\right).
\]
For fixed positive definite $\Sigma$ and $S$, let
\[
  c^\star(\Sigma,S) :=\argmin_{c>0}D(\Sigma,cS)=\Tr(\Sigma S^{-1})/d
\]
be the best scalar multiple of $S$.

\begin{lemma}[Consistency]\label{lem:consistency-t4}
Suppose the factor-gradient moments~\eqref{eq:observable-moments-t4} hold with
equality, with the model factors normalized so that
$\norm{\diag P^\star}_\infty=\norm{\diag Q^\star}_\infty=1$. For any diagonal
$\widehat P\succ0$ and $\widehat Q\succ0$, the
fits~\eqref{eq:factor-fit-t4} with $\widehat C_A=B^\top\widehat PB$ and
$\widehat C_B=A\widehat QA^\top$ return
\[
  \hat{q}(\widehat P)=\kappa\,c^\star\big(C^\star_A,\widehat C_A\big)\,\diag(Q^\star),
  \qquad
  \hat{p}(\widehat Q)=\kappa\,c^\star\big(C^\star_B,\widehat C_B\big)\,\diag(P^\star),
\]
so $\hat{q}/\norm{\hat{q}}_\infty=\diag(Q^\star)$ and
$\hat{p}/\norm{\hat{p}}_\infty=\diag(P^\star)$.
\end{lemma}
\begin{proof}
Under~\eqref{eq:observable-moments-t4}, $\Sigma_A$ is
block diagonal with $j$th block $\kappa\,q^\star_j\,C^\star_A$, and the
candidate $\diag(q)\otimes\widehat C_A$ of~\eqref{eq:factor-fit-t4} at a
point $q>0$ is block diagonal with $j$th block $q_j\,\widehat C_A$. Traces
and log-determinants add across blocks, so the objective is
$\sum_jD(\kappa q^\star_jC^\star_A,\,q_j\widehat C_A)$, one term per block,
and the $j$th term is minimized at its best scalar multiple,
\[
\hat{q}_j=c^\star(\kappa q^\star_jC^\star_A,\widehat C_A)
=\kappa\,q^\star_j\,c^\star(C^\star_A,\widehat C_A).
\]
The factor $\kappa\,c^\star(C^\star_A,\widehat C_A)$ is the same for every
$j$, and $\norm{\diag Q^\star}_\infty=1$, so the normalization removes it.
The claim for $\hat{p}$ follows with the roles of the two factors exchanged.
\end{proof}

\paragraph{Online estimator.} The fit~\eqref{eq:factor-fit-t4} averages past
factor gradients with weights $w_{t,s}$. \methodname{} implements this
average online. At step $t$ it computes one fitted term using
$C_{A,t}=B_t^\top P_tB_t$ and $C_{B,t}=A_tQ_tA_t^\top$, then folds that term
into an EMA,
\begin{equation}\label{eq:ours-klcoupling-t4}
\begin{aligned}
  \hat{q}_t(P_t) &= \frac1r\diag\!\left(G_{A,t}^\top C_{A,t}^{-1}G_{A,t}\right),
  &
  \hat{p}_t(Q_t) &= \frac1r\diag\!\left(G_{B,t} C_{B,t}^{-1}G_{B,t}^\top\right),
  \\
  q_{t+1} &= \beta_2q_t+(1-\beta_2)\hat{q}_t(P_t),
  &
  p_{t+1} &= \beta_2p_t+(1-\beta_2)\hat{p}_t(Q_t),
  \\
  Q_{t+1} &= \diag\!\bigl(q_{t+1}/\norm{q_{t+1}}_\infty\bigr),
  &
  P_{t+1} &= \diag\!\bigl(p_{t+1}/\norm{p_{t+1}}_\infty\bigr).
\end{aligned}
\end{equation}
Unrolling the EMA, up to the initialization tail, gives
\[
\begin{aligned}
  q_{t+1}
  &=
  \sum_{s\le t}w_{t,s}\hat q_s(P_s)
  \approx
  \frac1r\diag\!\left(
    \sum_{s\le t}w_{t,s}\,G_{A,s}^\top C_{A,t}^{-1}G_{A,s}
  \right)
  =
  \hat q(P_t),\\
  p_{t+1}
  &=
  \sum_{s\le t}w_{t,s}\hat p_s(Q_s)
  \approx
  \frac1r\diag\!\left(
    \sum_{s\le t}w_{t,s}\,G_{B,s}C_{B,t}^{-1}G_{B,s}^\top
  \right)
  =
  \hat p(Q_t).
\end{aligned}
\]
The exact recurrence is the EMA over the per-step fits
$\hat q_s(P_s),\hat p_s(Q_s)$. Reading this EMA as the weighted fit at the
current matrices replaces older $C_{A,s},C_{B,s}$ terms by
$C_{A,t},C_{B,t}$, using the same slow-change approximation as above. Under
the displayed moment equations, \Cref{lem:consistency-t4} then gives
\[
  Q_{t+1}
  =
  \diag\!\bigl(q_{t+1}/\norm{q_{t+1}}_\infty\bigr)
  \approx Q^\star_t,
  \qquad
  P_{t+1}
  =
  \diag\!\bigl(p_{t+1}/\norm{p_{t+1}}_\infty\bigr)
  \approx P^\star_t .
\]

\paragraph{Comparison with KL-Shampoo.} KL-Shampoo~\cite{klshampoo} also
fits Kronecker factors with the log-determinant divergence, and its moving
average has the same shape as~\eqref{eq:ours-klcoupling-t4}. The fit here
differs in three ways.
\begin{itemize}[leftmargin=1.5em,itemsep=1pt,topsep=2pt]
  \item \emph{No distributional assumptions.} KL-Shampoo fits the second
  moment of a stationary gradient distribution. Here $\EMA_{s \leq t}$ is a
  weighted sum over the realized gradients, not an expectation. The
  consistency claim is conditional on the displayed moment equations, not
  on a distributional limit.
  \item \emph{The fitted moments come from $G_A$ and $G_B$.} KL-Shampoo
  fits Kronecker factors to the second moment of the weight gradient $G$.
  Here $G$ is never formed. Each diagonal factor is fit from the second
  moment of its corresponding LoRA factor gradient, and the two fits are
  coupled through $C_A$ and $C_B$.
  \item \emph{Fixed-moment statement.} KL-Shampoo
  justifies its moving average as a stochastic proximal-gradient step on
  the divergence objective, correct in the limit. Here
  \Cref{lem:consistency-t4} is an algebraic statement for fixed moments.
  Conditional on the displayed moment equations, using an inexact estimate
  of the other factor only rescales the fitted vector, which the normalization
  removes, and the EMA is exact over the per-step fits.
\end{itemize}

\paragraph{Recovering the update of \Cref{alg:ours}.} Substituting the
fitted model $\kappa_t(Q_t\otimes P_t)$ for the second moment and the
look-ahead momentum $\widehat M_t$ for the batch gradient gives the
preconditioned spectral LMO~\eqref{eq:polorafull}
(\Cref{lem:pseudogradient-spectral}). The direction and magnitude steps of
\Cref{sec:sample-loss-lmo}, with the factor momenta standing in for the two
dense momentum projections, give the update~\eqref{eq:polora-update}. In
\Cref{alg:ours}, the curvature factors update after the direction step, so
step $t$ uses the estimate from step $t-1$. This avoids forming and
inverting $C_A$ and $C_B$ twice per optimizer step.

\subsection{Normalization Placement}
\label{sec:scale-placement-t4}
\methodname{} normalizes the accumulated vectors $q_{t+1},p_{t+1}$ when it
forms $Q_{t+1},P_{t+1}$ in~\eqref{eq:ours-klcoupling-t4}. It does not
normalize each estimate $\hat{q}_t,\hat{p}_t$ before adding it to the EMA.
The update is invariant to the overall scale of $P$ and $Q$, but the EMA
that estimates them is not.

\paragraph{Rescaling does not change the update.} The update depends on
$P$ and $Q$ only through their direction. Under a rescaling $Q\mapsto cQ$ with
$c>0$, both $Q^{-1/2}$ and $C_B^{-1/2}=(AQA^\top)^{-1/2}$ pick up a factor
$c^{-1/2}$. Inside the matrix sign this scalar has no effect, since the sign
is unchanged by positive rescaling. Outside it, the scalar multiplies $D_A$
and $D_B$ by $c^{-1/2}$, which the division by $\norm{D_A}_2$ and
$\norm{D_B}_2$ in the magnitude step cancels. The factor $P$ enters the same
way. The update therefore sees the direction of the curvature factors, but
not the magnitude $\kappa_t$ of the
model~\eqref{e-diag-kron-second-mom-model-t4}. Normalizing the accumulated
vectors therefore leaves it unchanged.

\paragraph{Rescaling changes the average.} The coupled fit makes each
estimate at step $t$ scale inversely with the curvature factor used to compute it,
\[
  \hat{q}_t(aP_t)=a^{-1}\,\hat{q}_t(P_t),
  \qquad
  \hat{p}_t(aQ_t)=a^{-1}\,\hat{p}_t(Q_t)
  \qquad(a>0),
\]
since $C_A=B^\top PB$ is linear in $P$, and likewise $C_B$ in $Q$.
Normalizing after averaging gives every factor $P_t,Q_t$ the same scale,
$\norm{\diag P_t}_\infty=\norm{\diag Q_t}_\infty=1$, so the estimates enter
the EMA on a common scale. If the factor used at step $s$
carried a scale $a_s$ that varied with $s$, the identity above would scale
$\hat{q}_s$ by $1/a_s$, so it would enter with effective weight
$w_{t,s}/a_s$ rather than the intended weight $w_{t,s}$ in
\Cref{sec:weighted-estimation-t4}. Normalizing the estimates before
averaging would also change the estimate, since normalization does not
commute with the mean. The same average-then-normalize order appears in
normalized SGD~\cite{cutkosky_nigt}, where momentum precedes normalization,
and in Muon~\cite{denoise_ortho}, where momentum precedes orthogonalization.
